% =========================================================================
% SENDA2 y el acomodo 3D — guía educativa 2 (repositorio LATTIMEX, versión 1.1.0)
% SPDX-License-Identifier: CC-BY-4.0
% Compilación: pdflatex guia_senda2 && bibtex guia_senda2 && pdflatex guia_senda2 && pdflatex guia_senda2
% Datos de las figuras y macros de cifras en datos2/ (generados por el autor desde los datos crudos).
% Esta guía no reemplaza a guia_senda.tex (versión 1.0.0), que se conserva sin cambios.
% =========================================================================
\documentclass[11pt,a4paper]{article}

\usepackage[T1]{fontenc}
\usepackage{lmodern}
\usepackage[spanish,es-noshorthands,es-nodecimaldot,es-tabla]{babel}
\usepackage[margin=2.4cm]{geometry}
\usepackage{microtype}
\usepackage{amsmath,amssymb}
\usepackage{booktabs}
\usepackage{array}
\usepackage{graphicx}
\usepackage{float}
\usepackage[font=small,labelfont=bf]{caption}
\usepackage{enumitem}
\usepackage[dvipsnames]{xcolor}
\usepackage{tikz}
\usetikzlibrary{arrows.meta,positioning,calc,shapes.geometric}
\usepackage{pgfplots}
\pgfplotsset{compat=1.18}
\usepackage[most]{tcolorbox}
\usepackage{algorithm}
\usepackage{algpseudocode}
\usepackage{url}
\usepackage[colorlinks=true,linkcolor=MidnightBlue,citecolor=OliveGreen,urlcolor=MidnightBlue]{hyperref}

\floatname{algorithm}{Algoritmo}
\algrenewcommand\algorithmicrequire{\textbf{Entrada:}}
\algrenewcommand\algorithmicensure{\textbf{Salida:}}
\algrenewcommand\algorithmicwhile{\textbf{mientras}}
\algrenewcommand\algorithmicdo{\textbf{hacer}}
\algrenewcommand\algorithmicif{\textbf{si}}
\algrenewcommand\algorithmicthen{\textbf{entonces}}
\algrenewcommand\algorithmicelse{\textbf{si no}}
\algrenewcommand\algorithmicfor{\textbf{para}}
\algrenewcommand\algorithmicend{\textbf{fin}}
\algrenewcommand\algorithmicreturn{\textbf{devolver}}
\algrenewtext{EndWhile}{\textbf{fin mientras}}
\algrenewtext{EndIf}{\textbf{fin si}}
\algrenewtext{EndFor}{\textbf{fin para}}

% --- Cajas didácticas (las mismas de la guía 1) ---------------------------
\definecolor{sendaink}{RGB}{19,34,31}
\definecolor{sendared}{RGB}{210,75,50}
\definecolor{sendagreen}{RGB}{11,131,96}
\newtcolorbox{ideaclave}[1][]{enhanced,breakable,colback=sendagreen!6,colframe=sendagreen!70!black,
  fonttitle=\bfseries,title={Idea clave#1},left=6pt,right=6pt,top=4pt,bottom=4pt}
\newtcolorbox{ejemplo}[1][]{enhanced,breakable,colback=MidnightBlue!4,colframe=MidnightBlue!70,
  fonttitle=\bfseries,title={Ejemplo resuelto#1},left=6pt,right=6pt,top=4pt,bottom=4pt}
\newtcolorbox{leccion}[1][]{enhanced,breakable,colback=sendared!6,colframe=sendared!85!black,
  fonttitle=\bfseries,title={Lección metodológica#1},left=6pt,right=6pt,top=4pt,bottom=4pt}
\newtcolorbox{reproducir}[1][]{enhanced,breakable,colback=gray!6,colframe=gray!60!black,
  fonttitle=\bfseries,title={Para reproducir#1},left=6pt,right=6pt,top=4pt,bottom=4pt}

\newcommand{\senda}{SENDA}
\newcommand{\sendados}{SENDA2}
\newcommand{\pp}{\,pp}
\newcommand{\tsl}[1]{\texttt{\small #1}}

% generado por tools/datos_guia2.py
\newcommand{\FiloABE}{0.053}
\newcommand{\FiloX}{0.72}
\newcommand{\HgsABE}{0.006}
\newcommand{\HgsX}{1.40}
\newcommand{\NABEtres}{55}
\newcommand{\ProdInvalidosMotor}{9}
\newcommand{\ProdRechazoMotor}{70.8}
\newcommand{\ProdRechazoPtresL}{14.6}
\newcommand{\ProdRechazoPtresLq}{3.8}
\newcommand{\SendaDosABE}{0.074}
\newcommand{\SendaDosGanaHgsX}{7}
\newcommand{\SendaDosX}{1.74}

% generado por tools/macros_motor.py
\newcommand{\AochoExGap}{+26.4\,\%}
\newcommand{\AochoExValid}{145/285}
\newcommand{\AochoProdLoValid}{pendiente}
\newcommand{\AochoProdRelaxed}{pendiente}
\newcommand{\AochoProdValid}{pendiente}
\newcommand{\AochoUpGap}{pendiente}
\newcommand{\AochoUpGapEasy}{pendiente}
\newcommand{\AochoUpGapHard}{pendiente}
\newcommand{\AochoUpValid}{pendiente}
\newcommand{\HashAocho}{pendiente}
\newcommand{\HashNivelUno}{1944dc5be91d}
\newcommand{\HashNivelUnoP}{e2bde2f206bd}
\newcommand{\HashRondaDos}{8dfb1d575211}
\newcommand{\HashSdosRtres}{ab706daaaabf}
\newcommand{\ReproABE}{0.086}
\newcommand{\ReproX}{1.77}
\newcommand{\ReproXParalelo}{1.96}
\newcommand{\VerifConcuerda}{2\,436}


\title{\textbf{SENDA2 y el acomodo tridimensional:\\ del ruteo clásico al problema integrado, medido con honestidad}\\[0.4em]
\large Segunda guía educativa del motor LATTIMEX (versión 1.1.0): qué cambió, cómo se midió\\ y cómo probarlo en su computadora}
\author{Erick Martínez\\
\small LATTIMEX · investigación independiente · Querétaro, México\\
\small \texttt{contacto@lattimex.com} · \texttt{https://lattimex.com}}
\date{Octubre de 2026}

\hypersetup{
  pdftitle={SENDA2 y el acomodo tridimensional: guía educativa 2 del motor LATTIMEX},
  pdfauthor={Erick Martínez (LATTIMEX)},
  pdfsubject={Ruteo de vehículos con carga tridimensional (3L-CVRP)},
  pdfkeywords={CVRP, 3L-CVRP, ALNS, acomodo 3D, orden de descarga, evaluación reproducible}
}

\makeatletter
\def\ps@licenciaguia{%
  \let\@oddhead\@empty
  \let\@evenhead\@empty
  \def\@oddfoot{\vtop{\offinterlineskip%
    \hbox to\textwidth{\hfil\thepage\hfil}%
    \vskip5pt
    \hbox to\textwidth{\hfil\scriptsize\textcopyright{} 2026 Erick Martínez (LATTIMEX).
      CC BY 4.0: \url{https://creativecommons.org/licenses/by/4.0/}\hfil}}}%
  \let\@evenfoot\@oddfoot
}
\makeatother

\begin{document}
\maketitle
\thispagestyle{licenciaguia}

\begin{abstract}
\noindent\textbf{En pocas palabras.} La primera guía abrió \senda, el núcleo de ruteo de LATTIMEX, y
lo midió en el problema clásico. Esta segunda guía cuenta el avance: un núcleo mejor (\sendados), un
acomodador tridimensional nuevo (P3L) y su integración en el motor 1.1.0, que es de código abierto
(MIT). Igual que la primera, explica cada pieza desde cero y dedica tanto espacio a \emph{cómo} se
midió como a los resultados, incluidos los errores propios que la medición destapó.

\smallskip
\noindent \sendados{} agrega a \senda{} tres cambios en su búsqueda adaptativa de vecindarios grandes:
una reparación con caché, el descarte de reparaciones que no caben y un tope de 0.02\,s al pulido de
cada nuevo mejor. En una misma sesión, con el mismo reloj y verificación independiente, queda a
+0.33\pp{} de HGS-CVRP en 41 instancias X a 4\,s (\SendaDosX\,\% frente a \HgsX\,\%; FILO2,
\FiloX\,\%) y a +0.07\pp{} en A/B/E a 2\,s con flota fija; la primera guía reportaba una distancia de
unos 2.6\pp{} en X. Para la carga se construyó un arnés contra los óptimos publicados de Tamke et al.\
(2025) y el validador de Krebs y Ehmke (2023). Con él se descubrió que el empacador de la versión 1.0.0
rechazaba entre 34\,\% y 61\,\% de rutas que sí caben y que permitía una caja posterior encima de una
anterior. P3L baja ese rechazo a 24--36\,\% y, con el verificador estricto, el motor ya no entrega
acomodos que violen el orden de descarga. Un ruteo que revisa la carga en cada paso (R3L) logra 285 de
285 soluciones válidas en las 19 instancias de Gendreau, a +6.9\,\% de distancia de las soluciones de
Tamke (que corrió hasta 8 horas por instancia). R3L no escala a 300 entregas; el documento explica por
qué y qué queda pendiente.
\end{abstract}

\noindent\textbf{Palabras clave:} ruteo de vehículos; CVRP; 3L-CVRP; ALNS; acomodo tridimensional;
orden de descarga (LIFO); evaluación reproducible; docencia en optimización.

\tableofcontents
\newpage

% =========================================================================
\section*{Cómo leer este documento}
\addcontentsline{toc}{section}{Cómo leer este documento}

Esta guía continúa a la primera (\emph{SENDA: ruteo de vehículos con búsqueda de trayectoria única,
explicado y medido con honestidad}, 30 de septiembre de 2026, \tsl{paper/guia\_senda.pdf}). No la
reemplaza: la primera sigue siendo la referencia para entender \senda{} pieza por pieza (construcción,
búsqueda local granular, ALNS, movimientos de segmento, portafolio) y para las bases de una
comparación justa. Aquí se supone ese conocimiento y se explica lo nuevo.

\begin{itemize}[leftmargin=1.5em]
\item La Sección~\ref{sec:sendados} presenta \sendados{}, cómo se llegó a él y cómo se compara con
HGS-CVRP y FILO2 en el problema clásico.
\item La Sección~\ref{sec:carga} introduce el problema con carga tridimensional, explica cómo medir
un acomodador sin engañarse y presenta P3L.
\item La Sección~\ref{sec:integracion} une ruteo y acomodo: por qué un buen acomodador no basta, qué
es R3L y qué logra el motor 1.1.0.
\item La Sección~\ref{sec:pruebalo} explica cómo usar el motor completo o cada pieza por separado.
\end{itemize}

Los recuadros tienen las mismas cuatro funciones que en la primera guía:

\begin{ideaclave}
Resume el concepto que no debe olvidarse.
\end{ideaclave}
\begin{ejemplo}
Resuelve un cálculo pequeño, paso a paso y con números concretos.
\end{ejemplo}
\begin{leccion}
Describe un error posible, o uno que realmente ocurrió en este proyecto, y cómo evitarlo.
\end{leccion}
\begin{reproducir}
Indica las huellas SHA-256 de los artefactos que respaldan una cifra y cómo repetir la medición.
\end{reproducir}

\paragraph{Sobre esta edición.} Acompaña a la versión 1.1.0 del repositorio público de LATTIMEX. El
repositorio contiene el motor y esta guía; no contiene los programas de comparación (HGS-CVRP, FILO2,
el validador de Krebs, el modelo exacto con CP-SAT) ni los scripts de las campañas. Se describen aquí
para que la comparación sea transparente; los datos crudos, contratos y bitácoras se conservan en el
archivo del autor y cada cifra cita la huella del artefacto que la respalda.

% =========================================================================
\section{Introducción}
\label{sec:intro}

La primera guía terminaba con tres tareas pendientes: dar más tiempo efectivo a la búsqueda en
instancias grandes, medir el núcleo en el problema integrado con carga tridimensional y ampliar el
banco de pruebas. Esta guía cuenta lo que pasó con las dos primeras.

La respuesta corta tiene tres partes:

\begin{enumerate}[leftmargin=1.5em]
\item \textbf{\sendados.} Tres cambios pequeños al ALNS de \senda{} cierran la mayor parte de la
distancia con HGS-CVRP en las instancias X: de unos 2.6 a 0.33 puntos de gap a 4\,s. El camino incluyó
un error de medición propio y una hipótesis (la semilla \emph{k-medoids}) que la ablación descartó.
\item \textbf{P3L y el verificador estricto.} Medir el empacador de la versión 1.0.0 contra rutas que
\emph{se sabe} que caben reveló que rechazaba la mitad de ellas y que su regla de orden de descarga
tenía un hueco. P3L es un acomodador nuevo con la semántica estándar; el verificador estricto garantiza
que el motor no entregue un acomodo inválido.
\item \textbf{La integración.} Un buen acomodador no basta: si el ruteo no sabe qué cabe, propone rutas
imposibles. R3L, un ruteo que consulta al acomodador en cada paso, resuelve la validez en instancias
pequeñas; a escala de producto el motor usa P3L dentro del ruteo de \sendados{} y avisa cuando no logra
respetar el orden de descarga.
\end{enumerate}

\begin{ideaclave}
Cada avance de esta guía empezó por una medición que contradijo lo que se creía. La versión 1.0.0
no era peor de lo anunciado en ruteo, pero sí en carga: con orden de descarga entregaba planes que lo
violaban. Medir contra óptimos conocidos y con un juez independiente fue lo que lo hizo visible.
\end{ideaclave}

% =========================================================================
\section{\sendados: el núcleo de ruteo}
\label{sec:sendados}

\subsection{Qué cambia respecto de \senda}

\senda{} (primera guía, Secciones 3 a 7) sigue una sola trayectoria por corrida: construye, mejora con
búsqueda local granular, reestructura con un ALNS que destruye parte de la solución y la repara, pule
cada nuevo mejor y reparte el tiempo entre dos o tres configuraciones. \sendados{} conserva todo eso y
cambia tres cosas dentro del ALNS:

\begin{enumerate}[leftmargin=1.5em,label=(\alph*)]
\item \textbf{Reparación regret con caché.} Al reinsertar los clientes quitados, la reparación de
\senda{} recalcula, para cada cliente y en cada paso, su mejor posición en todas las rutas. \sendados{}
guarda una sola inserción candidata por ruta para cada cliente (la más barata entre las posiciones
contiguas a sus vecinos granulares) y, después de cada inserción, sólo recalcula la ruta que cambió.
Cada reparación cuesta mucho menos y el ALNS hace más iteraciones en el mismo tiempo.
\item \textbf{Descarte de reparaciones fallidas.} Con flota fija y sin penalización, si la reparación
no logra meter a todos los clientes sin exceder la capacidad, \sendados{} descarta el candidato (y
castiga los pesos de los operadores que lo produjeron) en lugar de aceptarlo infactible.
\item \textbf{Tope de pulido.} Cada vez que aparece un nuevo mejor, \senda{} lo pule con búsqueda
local hasta que no mejora. \sendados{} limita ese pulido a $\min(0.02\text{\,s},\,0.20\cdot t_{\text{rest}})$,
donde $t_{\text{rest}}$ es el tiempo que queda.
\end{enumerate}

\begin{ejemplo}[: cuánto tiempo libera el tope de pulido]
En X a 4\,s, cada solver dispone del 80\,\% del reloj: 3.2\,s. En la medición previa a \sendados{}, el
pulido de cada nuevo mejor consumía en promedio alrededor de 1.9 de esos 3.2\,s, casi el 60\,\%. Con el
tope de 0.02\,s, un ensayo que encuentra 40 nuevos mejores gasta a lo sumo $40 \times 0.02 = 0.8$\,s en
pulir, y el resto vuelve al ALNS. En la confirmación (Sección~\ref{sec:confirmacion}), las iteraciones
medias en X pasaron de 654 a 1\,402 por ensayo y el gap bajó 0.59 puntos.
\end{ejemplo}

\subsection{Cómo se llegó: dos errores útiles}

\begin{leccion}[: el control degradado]
El primer candidato, llamado \emph{KM-R1}, sembraba la solución inicial con $k$-medoides y destruía
por medoides. Su primera comparación contra \senda{} mostró una ventaja grande. Al repetirla se
descubrió que el control (\senda) había corrido a la mitad de sus iteraciones normales: otra sesión,
otra carga de la máquina. Repetida con los tres brazos en la misma sesión y con un \emph{sentinela de
velocidad} (las iteraciones del control deben coincidir con las de una corrida de referencia), la
ventaja seguía siendo real pero menor. Regla: \textbf{todos los brazos en la misma sesión, con orden
rotado y un sentinela de velocidad}; nunca contra un número de otra sesión.
\end{leccion}

\begin{leccion}[: el nombre que la ablación quitó]
Con la ventaja confirmada, era tentador atribuirla a los $k$-medoides. La ablación (quitar una pieza
a la vez, en la misma sesión) dijo otra cosa: quitar la semilla o la destrucción por medoides no
cambiaba nada; quitar la reparación con caché costaba 0.87 puntos en X, y el tope de pulido explicaba
la mayor parte del resto. El candidato se simplificó a esas piezas y perdió el nombre ``SENDA-KM'':
se llama \sendados. Además, el tope de 0.02\,s se eligió mirando los datos de una ronda, así que su
cifra no podía publicarse sin repetirla con semillas nuevas (Sección~\ref{sec:confirmacion}).
\end{leccion}

\subsection{Resultados frente a HGS-CVRP y FILO2}
\label{sec:confirmacion}

\paragraph{Diseño.} Igual que en la primera guía: 60 instancias A/B/E con flota fija a 2\,s y 41
instancias X con flota libre a 4\,s; el 80\,\% del reloj como presupuesto interno para todos los
solvers; un hilo; distancias redondeadas al entero (convención de CVRPLIB); verificación
independiente de cobertura, capacidad, flota y costo. HGS-CVRP~\cite{vidal2022hgs} y
FILO2~\cite{accorsi2023} se compilaron de sus fuentes oficiales. La tabla principal es una sesión con
los tres solvers y semillas nunca usadas (3000--3900); la confirmación H1-C, con otras semillas nuevas
(2000--2900), repite la comparación con HGS.

\begin{table}[H]
\centering
\caption{Gap medio contra la mejor solución conocida (BKS). Misma sesión, mismo reloj. En A/B/E se
promedian las \NABEtres{} instancias donde los tres solvers respetaron la flota en todas sus corridas
(FILO2 no tiene opción de flota fija y usó una ruta de más en 5 instancias).}
\label{tab:cvrp}
\begin{tabular}{lrrr}
\toprule
& \sendados & HGS-CVRP & FILO2 \\
\midrule
X, 41 instancias, 4\,s, 5 semillas & \SendaDosX\,\% & \HgsX\,\% & \textbf{\FiloX\,\%} \\
A/B/E, \NABEtres{} instancias, 2\,s, 10 semillas & \SendaDosABE\,\% & \textbf{\HgsABE\,\%} & \FiloABE\,\% \\
\bottomrule
\end{tabular}
\end{table}

\paragraph{Comparaciones pareadas.} Por instancia, $\Delta$ = gap de \sendados{} menos gap del otro
solver (IC95 por bootstrap; prueba de Wilcoxon):
\begin{itemize}[leftmargin=1.5em]
\item \textbf{Contra HGS, en X:} $+0.33$\pp{} [0.21, 0.45], $p = 3\cdot10^{-6}$; \sendados{} es mejor
en \SendaDosGanaHgsX{} de 41 instancias. La confirmación H1-C da $+0.34$\pp{} con \sendados{} y
$+0.28$\pp{} [0.13, 0.42] con una variante que conservaba las piezas descartadas.
\item \textbf{Contra HGS, en A/B/E:} $+0.068$\pp{} [0.033, 0.111]; empata en 36 instancias y HGS es
mejor en 19.
\item \textbf{Contra FILO2:} en X, $+1.02$\pp{} [0.85, 1.17], y FILO2 gana en 40 de 41; en A/B/E no
se detectó una diferencia clara ($+0.021$\pp{}, IC que cruza cero).
\end{itemize}

\begin{figure}[H]
\centering
\begin{tikzpicture}
\begin{axis}[width=0.92\textwidth,height=6.2cm,xlabel={clientes ($n$)},ylabel={gap medio (\%)},
  legend style={at={(0.5,1.03)},anchor=south,legend columns=3,draw=none},legend cell align=left,grid=major,grid style={gray!20},ymin=0]
\addplot[only marks,mark=*,mark size=1.5pt,sendared] table[x=n,y=senda2]{datos2/x_senda2.dat};\addlegendentry{\sendados}
\addplot[only marks,mark=square*,mark size=1.5pt,MidnightBlue] table[x=n,y=hgs]{datos2/x_senda2.dat};\addlegendentry{HGS-CVRP}
\addplot[only marks,mark=triangle*,mark size=1.8pt,sendagreen] table[x=n,y=filo2]{datos2/x_senda2.dat};\addlegendentry{FILO2}
\end{axis}
\end{tikzpicture}
\caption{Gap medio por instancia X a 4\,s (5 semillas por solver, misma sesión). La ventaja de FILO2
crece con el tamaño de la instancia.}
\label{fig:x}
\end{figure}

\paragraph{Contra \senda.} La primera guía reportaba, en otra máquina, 4.119\,\% en X frente a
1.450\,\% de HGS. Para no mezclar sesiones, la comparación honesta usa a HGS como ancla común: en una
sesión de octubre, \senda{} quedaba a $+2.64$\pp{} de HGS en X (4.05\,\% frente a 1.41\,\%);
\sendados{} queda a $+0.33$\pp. En A/B/E, de $+0.094$ a $+0.068$\pp.

\begin{ideaclave}
La afirmación honesta no es ``estado del arte'': FILO2 es claramente más fuerte en instancias grandes
con reloj corto, y HGS sigue siendo mejor en promedio. \sendados{} es un núcleo de unas 2\,200 líneas,
de una sola trayectoria, que respeta la flota fija y queda cerca de HGS. Su valor está en ser simple,
auditable y medido con cuidado.
\end{ideaclave}

\begin{reproducir}
Sesión con los tres solvers: \tsl{raw\_trials.jsonl} \tsl{ef823740c317\ldots}, contrato
\tsl{906e37b0ab29\ldots}. Confirmación H1-C: \tsl{2ac424ad3411\ldots}. Con el repositorio:
\tsl{python -m lattimex cvrp INSTANCIA.vrp --time 4} (X) o \tsl{--time 2 --fleet fixed} (A/B/E)
reparte el tiempo y elige configuraciones igual que en la medición. Repetido sobre el mismo banco con
el binario publicado, una corrida a la vez, dio \ReproX\,\% en X (41 instancias, 1 semilla) y
\ReproABE\,\% en A/B/E (60 instancias, 2 semillas), dentro de la variación entre semillas.
\textbf{Corra una instancia a la vez}: con seis procesos en paralelo en la misma máquina el gap en X
subió a \ReproXParalelo\,\%, porque cada proceso hizo menos iteraciones.
\end{reproducir}

\subsection{\sendados{} dentro del motor}
\label{sec:motor-ruteo}

El motor no llama al núcleo igual que el banco: reparte el presupuesto entre ruteo, acomodo y
reoptimización, usa flota exacta y corre un portafolio de configuraciones en dos hilos. Medido así
(107 instancias, 3 semillas, misma sesión), \sendados{} contra \senda{}:

\begin{itemize}[leftmargin=1.5em]
\item A/B/E: $-0.036$\,\% de distancia [$-0.068$, $-0.009$], $p=0.02$.
\item X: el resultado depende de cuántos clientes atiende cada ruta (Spearman $\rho=-0.65$). Con 12 o
más clientes por ruta, $-0.89$\,\% (gana en 14 de 16); con menos de 6, $+2.99$\,\% (pierde en 9 de 12).
\end{itemize}

Las rutas reales de LATTIMEX son largas (de 30 a 75 entregas por unidad), así que el motor 1.1.0 usa
\sendados{} por omisión y vuelve al núcleo de 1.0.0 cuando hay menos de 6 entregas por unidad. La
regla se fijó con los mismos datos que la motivan; es una decisión de ingeniería, no una cifra
confirmada.

La misma ronda encontró un riesgo del constructor de flota exacta: cuando la carga ocupa entre 92.5\,\%
y 99.7\,\% de la capacidad total, no lograba repartirla en $k$ unidades y el motor rechazaba el pedido
(11 de 41 instancias X). El motor 1.1.0 agrega un respaldo que reparte aceptando sobrecarga y la
elimina con búsqueda local; resolvió las 11 en menos de medio segundo.

% =========================================================================
\section{La carga: medir y construir un acomodador}
\label{sec:carga}

\subsection{El problema con carga tridimensional}

En el 3L-CVRP~\cite{gendreau2006} cada cliente pide cajas con largo, ancho, alto, peso y una marca de
fragilidad, y cada ruta debe poder acomodarse dentro de la caja del vehículo. Se usan los ejes de
Gendreau: $x$ es el largo (la puerta está en $x=L$), $y$ el ancho y $z$ el alto; las cajas sólo rotan
en horizontal. Además de no salirse ni encimarse, un acomodo puede exigir:

\begin{itemize}[leftmargin=1.5em]
\item \textbf{Soporte:} al menos el 75\,\% de la base de cada caja descansa sobre el piso o sobre
cajas cuya cara superior está a esa altura.
\item \textbf{Fragilidad:} una caja no frágil no puede descansar sobre una frágil.
\item \textbf{Orden de descarga (LIFO):} al llegar con cada cliente, sus cajas salen sin mover las de
clientes posteriores.
\end{itemize}

Las cinco variantes estándar combinan estas reglas:

\begin{table}[H]
\centering
\caption{Variantes del 3L-CVRP (Gendreau et al., 2006).}
\label{tab:variantes}
\begin{tabular}{lccc}
\toprule
Variante & Soporte & Fragilidad & Orden de descarga \\
\midrule
loading-only & -- & -- & -- \\
no-support & -- & sí & sí \\
no-fragility & sí & -- & sí \\
no-lifo & sí & sí & -- \\
all-constraints & sí & sí & sí \\
\bottomrule
\end{tabular}
\end{table}

\paragraph{El orden de descarga, con precisión.} Sea $e$ una caja de un cliente que se visita
\emph{antes} que el de la caja $t$. El acomodo viola el orden si $t$ bloquea a $e$ de alguna de dos
formas (Figura~\ref{fig:lifo}): (i) $t$ queda \emph{delante} de $e$, más cerca de la puerta, y sus
proyecciones sobre el plano $y$--$z$ se tocan; o (ii) $t$ queda \emph{encima} de $e$ y sus proyecciones
sobre el plano $x$--$y$ se tocan. Ésta es la semántica del validador público de Krebs y
Ehmke~\cite{krebs2023,krebsvalidator}.

\begin{figure}[H]
\centering
\begin{tikzpicture}[scale=0.9]
  \begin{scope}
    \draw[thick] (0,0) rectangle (6,3); \node[below] at (3,0) {vista lateral ($x$--$z$)};
    \draw[very thick,sendagreen] (6,0) -- (6,3); \node[right,sendagreen] at (6,1.5) {puerta};
    \filldraw[fill=MidnightBlue!25] (1,0) rectangle (3,1.4); \node at (2,0.7) {$e$ (1.ª)};
    \filldraw[fill=sendared!30] (3.2,0) rectangle (5.4,1.4); \node at (4.3,0.7) {$t$ (2.ª)};
    \node[align=center,font=\small] at (3,2.4) {(i) $t$ delante de $e$: \textbf{viola}};
  \end{scope}
  \begin{scope}[xshift=7.6cm]
    \draw[thick] (0,0) rectangle (6,3); \node[below] at (3,0) {vista lateral ($x$--$z$)};
    \draw[very thick,sendagreen] (6,0) -- (6,3); \node[right,sendagreen] at (6,1.5) {puerta};
    \filldraw[fill=MidnightBlue!25] (2,0) rectangle (5,1.2); \node at (3.5,0.6) {$e$ (1.ª)};
    \filldraw[fill=sendared!30] (2.6,1.2) rectangle (4.4,2.2); \node at (3.5,1.7) {$t$ (2.ª)};
    \node[align=center,font=\small] at (3,2.6) {(ii) $t$ encima de $e$: \textbf{viola}};
  \end{scope}
\end{tikzpicture}
\caption{Las dos formas en que una caja $t$ de una entrega posterior bloquea a una caja $e$ de una
entrega anterior. El empacador de la versión 1.0.0 sólo prohibía la primera.}
\label{fig:lifo}
\end{figure}

\subsection{Cómo medir un acomodador sin engañarse}

Medir un acomodador parece fácil: se le dan rutas y se cuenta cuántas acomoda. La trampa es que no se
sabe cuántas \emph{deberían} caber. Si se le dan rutas holgadas, cualquier acomodador parece bueno; si
se le dan rutas imposibles, cualquiera parece malo.

\paragraph{Rutas que se sabe que caben.} Tamke et al.~\cite{tamke2025} resolvieron con un
\emph{branch-and-cut} exacto las 19 instancias clásicas de Gendreau con hasta 50 clientes, en las
cinco variantes, y publicaron sus soluciones con el acomodo de cada ruta~\cite{tamkerepo}. Cada una de
esas 648 rutas cabe: su acomodo existe y el validador lo aprueba. Si un acomodador rechaza una de ellas,
es un \textbf{falso rechazo}, sin ambigüedad. A esta prueba se le llama \emph{nivel 1}.

\paragraph{Un juez independiente.} Todo acomodo aceptado se revisa con el validador de Krebs, que no
forma parte del motor. Antes de usarlo como juez se comprobó que juzga bien: aprueba las 95
soluciones de Tamke (19 instancias $\times$ 5 variantes), rechaza las 76 soluciones de variantes laxas
evaluadas con todas las restricciones y rechaza los 19 casos con un solape forzado.

\begin{leccion}[: el 2.4\,\% que no medía lo que parecía]
En julio se había reportado que el empacador del motor rechazaba sólo el 2.4\,\% de las rutas. Esa
medición usaba rutas de otra heurística construidas con todas las restricciones y las probaba en la
variante más laxa: eran rutas holgadas. Con las rutas óptimas, que están apretadas, el rechazo era del
34 al 64\,\%. Una cifra de calidad sólo significa algo junto con el banco en que se midió.
\end{leccion}

\subsection{Lo que encontró la medición en la versión 1.0.0}

\begin{itemize}[leftmargin=1.5em]
\item \textbf{Falso rechazo alto:} entre 36\,\% (loading-only) y 64\,\% (all-constraints) de las rutas
óptimas, y crece con el tamaño de la ruta: con 9 a 12 cajas, 47 de 49 rechazadas. Con 1\,024
reinicios en lugar de los habituales, all-constraints sólo bajó de 64\,\% a 60\,\%: el límite era el
algoritmo de colocación, no el esfuerzo.
\item \textbf{Un hueco en el orden de descarga:} el empacador prohibía el caso (i) de la
Figura~\ref{fig:lifo} pero no el (ii). Krebs rechazó 27 acomodos que el motor daba por buenos.
\item \textbf{En el motor completo,} con todas las restricciones, 1 de 57 corridas fue válida: cuando
una ruta no cabía, el motor relajaba el orden de descarga en silencio.
\end{itemize}

\subsection{P3L}

P3L (\tsl{native/p3l.cpp}) acomoda una ruta en un contenedor con la semántica de la
Figura~\ref{fig:lifo}:

\begin{enumerate}[leftmargin=1.5em]
\item \textbf{Secuencia de carga.} Con orden de descarga, las cajas se cargan del último cliente al
primero; dentro de cada cliente, por volumen, área de base, altura o lado mayor.
\item \textbf{Colocación por gravedad.} Para cada caja se generan posiciones candidatas en $x$ y $y$
(el origen, el final de cada caja ya puesta y la pared opuesta) y cada una baja hasta descansar;
gana la primera según una regla lexicográfica (por ejemplo, menor $x$, luego menor $z$).
\item \textbf{Búsqueda en haz} sobre la secuencia: se conservan las mejores secuencias parciales según
el contacto con paredes y cajas.
\item \textbf{Búsqueda local iterada} sobre la secuencia de carga y las rotaciones, desde varios
arranques, hasta colocar todo o agotar el tope (50\,ms por ruta en el motor).
\end{enumerate}

\begin{leccion}[: posiciones que parecen bastar y no bastan]
Para acelerar, es común restringir las posiciones a \emph{patrones normales}: las sumas de longitudes
de cajas. Es correcto sin soporte ni orden de descarga, pero no con ellos. Un modelo exacto con
CP-SAT, usado como referencia durante el desarrollo, ``demostró'' imposibles 5 rutas que Tamke había
acomodado. La causa eran los patrones normales; con todas las posiciones enteras en $x$ y $y$ cuando
hay soporte, como hace Tamke, las 5 rutas cupieron.
\end{leccion}

\paragraph{Resultados en el nivel 1.} Las 13 instancias que no se usaron para ajustar P3L, con el
candidato congelado (Tabla~\ref{tab:nivel1}). La columna ``P3L + CP-SAT'' agrega el modelo exacto para
las rutas que P3L no acomoda; baja el rechazo casi a cero, pero no forma parte del motor: depende de
OR-Tools y no termina con rutas de 55 o más cajas.

\begin{table}[H]
\centering
\caption{Falso rechazo sobre las rutas óptimas de Tamke, 13 instancias de confirmación. Ningún acomodo
aceptado por P3L fue rechazado por el validador.}
\label{tab:nivel1}
\begin{tabular}{lrrr}
\toprule
Variante & Empacador 1.0.0 & \textbf{P3L} (motor 1.1.0) & P3L + CP-SAT \\
\midrule
loading-only & 34.1\,\% & 23.9\,\% & 1.1\,\% \\
no-support & 42.2\,\% & 33.3\,\% & 0.0\,\% \\
no-fragility & 57.6\,\% & 32.6\,\% & 0.0\,\% \\
no-lifo & 41.1\,\% & 35.6\,\% & 6.7\,\% \\
all-constraints & 61.3\,\% & 30.1\,\% & 0.0\,\% \\
\bottomrule
\end{tabular}
\end{table}

\paragraph{A escala de producto.} En 159 rutas de 20 a 97 cajas producidas por el motor en 12
instancias de Querétaro, con todas las restricciones, el empacador de 1.0.0 rechazó al menos el
\ProdRechazoMotor\,\% de las rutas que caben y dio \ProdInvalidosMotor{} acomodos inválidos; P3L con
50\,ms, al menos el \ProdRechazoPtresL\,\%, y con 500\,ms, al menos el \ProdRechazoPtresLq\,\%, sin
ningún inválido (``al menos'' porque aquí no hay óptimos: una ruta cuenta como factible si algún
método la acomodó). Sin orden de descarga, ambos acomodan todas.

\paragraph{El verificador estricto.} El motor 1.1.0 revisa todo acomodo con la semántica de la
Figura~\ref{fig:lifo} antes de aceptarlo, venga de P3L o del empacador de 1.0.0, que se conserva como
respaldo porque puede girar en vertical las cajas que lo permiten. Contra el validador de Krebs, el
verificador coincidió en \VerifConcuerda{} de \VerifConcuerda{} acomodos, incluidos 105 que Krebs
rechaza.

% =========================================================================
\section{La integración: ruteo que sabe lo que cabe}
\label{sec:integracion}

\subsection{Por qué un buen acomodador no basta}

El motor rutea con una capacidad de volumen efectiva, una fracción del volumen del vehículo, y acomoda
después. Si una ruta no cabe, una \emph{mediación} prueba unos pocos cambios: invertirla, rotarla o
mover un cliente a otra ruta. Al sustituir sólo el acomodador por P3L, las corridas válidas con todas
las restricciones subieron de 1 a 20 de 57, pero el gap de las válidas no cambió (+15 a +29\,\%) y el
modelo exacto demostró imposibles 1\,665 rutas que el ruteo había propuesto. El ruteo seguía
proponiendo rutas sin saber si cabían.

\subsection{R3L}

R3L (\tsl{lattimex/engine/r3l.py}) es un ruteo de destruir y reconstruir al estilo
SISR~\cite{christiaens2020} con recocido simulado~\cite{kirkpatrick1983}, en el que cada inserción
consulta al acomodador:

\begin{algorithm}[H]
\caption{R3L (esquema)}
\begin{algorithmic}[1]
\Require clientes, matriz de distancias, contenedor, variante, flota máxima $k$, tiempo $T$
\State $S \gets$ reconstruir desde cero; si hay rutas de \sendados, también desde ellas, y quedarse con la mejor
\While{quede tiempo}
  \State $S' \gets$ quitar de $S$ cadenas de clientes cercanos (SISR)
  \For{cada cliente quitado, en un orden al azar o por volumen o distancia}
    \State probar sus posiciones de la más barata a la más cara;
    \State quedarse con la primera cuya ruta \textbf{cabe} (P3L); abrir ruta nueva sólo si hay menos de $k$
  \EndFor
  \State aceptar $S'$ con el criterio de recocido simulado; guardar la mejor solución completa
\EndWhile
\end{algorithmic}
\end{algorithm}

Tres detalles lo hacen rápido. (1) La revisión se escalona: una pasada voraz de P3L para cada posición
candidata (0.3\,ms con rutas pequeñas) y la búsqueda en haz sólo para la posición más barata. (2) Un
caché guarda cada veredicto; sin orden de descarga la llave es el conjunto de clientes, porque el orden
no cambia el acomodo. (3) Con orden de descarga, si una ruta no cabe se prueba invertida: tiene la
misma distancia y otro orden de descarga.

\begin{ejemplo}[: la ruta invertida]
Una ruta visita a A y luego a B. A pide una caja no frágil que ocupa todo el piso; B, una frágil del
mismo tamaño. En el orden A, B, la caja de B (segunda entrega) debe cargarse primero, al fondo o
abajo; como ocupa todo el piso, queda abajo y la de A encima: una caja no frágil sobre una frágil,
prohibido. En el orden B, A, la de A va abajo y la frágil de B encima: permitido, y B sale primero sin
mover nada. La distancia es la misma; R3L sólo tiene que probar la ruta en ambos sentidos.
\end{ejemplo}

\subsection{Resultados contra Tamke}

Las 19 instancias de Gendreau, las cinco variantes y tres semillas (285 corridas por brazo), con flota
``hasta $k$'' ($k$ = vehículos disponibles, igual que Tamke). Una corrida es \emph{válida} si el
validador de Krebs aprueba ruteo y carga en la variante. El gap es la distancia contra la solución de
Tamke.

\begin{table}[H]
\centering
\caption{Motor completo contra Tamke. ``Fáciles'': instancias que Tamke resolvió al óptimo en menos de
1\,000\,s; en las demás, Tamke corrió hasta 8 horas y no cerró su propia cota.}
\label{tab:integracion}
\begin{tabular}{lrrrr}
\toprule
& válidas & gap medio & gap fáciles & gap demás \\
\midrule
Motor 1.0.0 (10\,s, flota exacta) & 81/285 & +26.4\,\% & -- & -- \\
\sendados{} + R3L, investigación (10\,s) & 283/285 & +8.7\,\% & +4.2\,\% & +13.5\,\% \\
\sendados{} + R3L, investigación (30\,s) & 285/285 & +6.9\,\% & +3.7\,\% & +10.2\,\% \\
\textbf{Motor 1.1.0}, flota exacta (10\,s) & \AochoExValid & \AochoExGap & -- & -- \\
\textbf{Motor 1.1.0}, \tsl{fleet\_mode=up\_to} (10\,s) & \AochoUpValid & \AochoUpGap & \AochoUpGapEasy & \AochoUpGapHard \\
\bottomrule
\end{tabular}
\end{table}

Con flota exacta (la opción por omisión), el motor 1.1.0 no usa R3L: rutea con \sendados{} y
acomoda con P3L. Casi duplica las corridas válidas del motor 1.0.0, pero el gap sigue alto, porque obliga
a usar todas las unidades disponibles y Tamke usa sólo las necesarias. Cuando relajó el orden de
descarga, el motor lo informó en la respuesta; el validador sólo rechazó corridas así.

\begin{leccion}[: una meta que no era realista]
La meta inicial era 100\,\% de corridas válidas y gap $\le 3$\,\% en 10\,s. Los tiempos de Tamke
muestran que no era razonable: en las instancias difíciles, su método exacto corrió 8 horas sin
demostrar el óptimo. La meta se ajustó después de ver los primeros resultados a $\le 3$\,\% en las
instancias fáciles, y se reporta aquí junto con la original. En las fáciles el promedio es de
+3.7\,\% a 30\,s (+3.3\,\% en las 13 instancias que no se usaron para ajustar R3L): todavía no se
cumple.
\end{leccion}

\subsection{A escala de producto}

R3L no escala: cada revisión vuelve a acomodar la ruta completa, y con rutas de 60 a 75 cajas una
revisión cuesta de 15 a 300\,ms, en lugar de los 0.3\,ms de las rutas de Gendreau. Con 300 entregas,
R3L no terminaba ni la solución inicial en el tiempo disponible. Por eso el motor 1.1.0 sólo lo usa
con flota ``hasta $k$'' en instancias pequeñas (hasta 60 entregas y 30 cajas por unidad).

Para instancias grandes, el motor 1.1.0 rutea con \sendados, acomoda con P3L en la variante pedida,
media con el orden de descarga estricto y, sólo si una ruta sigue sin caber, la acomoda sin orden de
descarga \textbf{y lo informa} en la respuesta (\tsl{lifo\_estricto}, \tsl{advertencias}). En las 12
instancias de producto con todas las restricciones (Tabla~\ref{tab:producto}), la versión 1.0.0
declaraba factibles todas las corridas y ninguna respetaba el orden de descarga.

\begin{table}[H]
\centering
\caption{12 instancias de producto (60 a 300 entregas, 1 a 10 unidades), flota exacta, 30\,s, 2
semillas. ``Con orden estricto'': el motor la declara factible, no relajó ninguna ruta y el validador
aprueba la carga.}
\label{tab:producto}
\begin{tabular}{lrr}
\toprule
& Motor 1.0.0 & \textbf{Motor 1.1.0} \\
\midrule
all-constraints: con orden estricto & 0/24 & \AochoProdValid \\
all-constraints: relajadas e \emph{informadas} & -- & \AochoProdRelaxed \\
loading-only: válidas & 24/24 & \AochoProdLoValid \\
\bottomrule
\end{tabular}
\end{table}

\begin{leccion}[: ``factible'' no es ``correcto'']
La versión 1.0.0 decía ``factible'' a planes que violaban el orden de descarga, porque lo relajaba sin
avisar para no partir rutas. Era una decisión de diseño razonable, pero invisible para quien leía la
respuesta. La 1.1.0 mantiene la posibilidad de relajar como último recurso y la hace visible: un plan
que obliga al operador a mover cajas no debe presentarse igual que uno que no.
\end{leccion}

% =========================================================================
\section{Discusión}
\label{sec:discusion}

\paragraph{Dónde está el valor.} En el problema clásico, \sendados{} acerca a LATTIMEX a HGS sin
volverlo más complejo. Pero dentro del motor, con carga tridimensional, la diferencia entre núcleos
de ruteo se pierde (la ronda 2 no detectó diferencia de distancia con acomodo 3D). La ganancia
grande para un operador está en la carga: que el plan respete el orden de descarga o diga cuándo no
puede.

\paragraph{Limitaciones.}
(i) El nivel 1 y el nivel 2 usan las 19 instancias de Gendreau con hasta 50 clientes; a escala de
producto no hay óptimos y las cifras son cotas.
(ii) P3L sólo gira cajas en horizontal; el empacador de 1.0.0 se conserva como respaldo para cajas que
pueden girar en vertical, sin una medición propia de ese caso.
(iii) La regla que elige entre \senda{} y \sendados{} se fijó con los mismos datos que la motivan.
(iv) R3L se ajustó con 6 instancias y se reporta también en las 13 restantes, pero las cifras de la
Tabla~\ref{tab:integracion} mezclan ambos grupos.
(v) Todas las mediciones son de una sola máquina.

\section{Conclusiones y trabajo futuro}

\sendados{} cierra la mayor parte de la distancia con HGS en el CVRP clásico con tres cambios
pequeños y medibles. P3L y el verificador estricto corrigen un defecto real de la versión 1.0.0 y
reducen a la mitad el rechazo de rutas que caben. R3L muestra que revisar la carga dentro del ruteo
resuelve la validez en instancias pequeñas.

\paragraph{Trabajo futuro.} (1) Un acomodo \emph{incremental}: al insertar un cliente como primera
entrega, sus cajas son las últimas en cargarse y podrían añadirse sin reacomodar la ruta; es la pieza
que falta para que R3L escale. (2) Un acomodador rápido más fuerte en rutas apretadas, para cerrar la
meta de $\le 3$\,\% en las instancias fáciles. (3) Localizar la búsqueda local de \sendados{} en los
vértices tocados, como FILO2, la diferencia de diseño más probable detrás de la brecha en X.

% =========================================================================
\section{Pruébelo en su computadora}
\label{sec:pruebalo}

El motor 1.1.0 está publicado con licencia MIT en \url{https://github.com/Erick-MF28/lattimex}; la
versión 1.0.0 sigue disponible en su etiqueta. Se puede usar completo o por piezas:

\begin{verbatim}
git clone https://github.com/Erick-MF28/lattimex.git
cd lattimex
pip install -r requirements.txt  # Windows sin compilador: pip install ziglang
python -m lattimex build         # compila los núcleos y los acomodadores
python -m lattimex probe         # huella determinista
python -m unittest discover -s tests/engine
\end{verbatim}

\paragraph{1. \sendados{} solo, en CVRPLIB.}
\begin{verbatim}
python -m lattimex cvrp X-n101-k25.vrp --time 4 --bks 27591
python -m lattimex cvrp A-n32-k5.vrp --time 2 --fleet fixed --sol A-n32-k5.sol
\end{verbatim}
\tsl{--router senda} usa el núcleo de la versión 1.0.0, para comparar ambos en su máquina.

\paragraph{2. El acomodador solo.} Un archivo JSON describe el contenedor y las cajas de una ruta,
con \tsl{stop} = posición de entrega (1 = primera):
\begin{verbatim}
{"cargo_cm": [300, 160, 160],
 "constraints": {"support": true, "fragility": true, "lifo": true},
 "items": [{"id": "A1", "dims_cm": [60, 40, 30], "stop": 1, "fragile": false},
           {"id": "B1", "dims_cm": [40, 30, 30], "stop": 2, "fragile": true}]}
\end{verbatim}
\tsl{python -m lattimex pack cajas.json} devuelve la posición de cada caja y si el acomodo cumple la
variante.

\paragraph{3. El motor completo.} \tsl{python -m lattimex serve} y el LATTIMEX Planner en
\url{https://lattimex.com/planner}, igual que en la versión 1.0.0: el Planner usa el motor 1.1.0
sin cambiar nada. Por API, el campo \tsl{fleet\_mode} con valor \tsl{up\_to} permite usar menos
unidades (y activa R3L en instancias pequeñas), y el campo \tsl{router} con valor \tsl{senda} fuerza
el núcleo de 1.0.0.

\paragraph{4. Repetir la medición de carga.} Las instancias y soluciones de Tamke~\cite{tamkerepo}
(MIT) y el validador de Krebs~\cite{krebsvalidator} (GPL-3.0) se obtienen de sus repositorios; no
forman parte de LATTIMEX. El validador lee las restricciones por número de línea: genere los cinco
archivos de variante a partir de su plantilla. Las soluciones de Tamke repiten el identificador de
tour, que hay que renumerar antes de validarlas.

\begin{reproducir}[: artefactos]
\small
\begin{tabular}{@{}p{5.2cm}p{9.3cm}@{}}
\sendados, HGS y FILO2 & \tsl{raw\_trials.jsonl} \tsl{ef823740c317\ldots}; contrato \tsl{906e37b0ab29\ldots} \\
Confirmación H1-C & \tsl{raw\_trials.jsonl} \tsl{2ac424ad3411\ldots} \\
\sendados{} en el motor (ronda 2) & \tsl{summary.json} \tsl{\HashRondaDos\ldots} \\
Nivel 1 (empacador 1.0.0 y P3L) & \tsl{level1.jsonl} \tsl{\HashNivelUno\ldots}; \tsl{level1\_hyb\_confirm.jsonl} \tsl{\HashNivelUnoP\ldots} \\
Integración (investigación) & \tsl{level2\_s2r3l\_summary.json} \tsl{\HashSdosRtres\ldots} \\
Motor 1.1.0 publicado & \tsl{resumen.json} \tsl{\HashAocho\ldots} \\
Figuras de este documento & datos en \tsl{paper/datos2/} de este repositorio \\
Artefactos completos & archivo del autor; se comparten a solicitud (\tsl{contacto@lattimex.com}) \\
\end{tabular}
\end{reproducir}

\bibliographystyle{plain}
\bibliography{referencias}

\end{document}
